\documentclass[12pt]{article}
\usepackage[utf8]{inputenc}
\usepackage[T1]{fontenc}
\usepackage{hyperref}
\title{Kolbila Language Resources for K-12 Curriculum Development}
\author{Compiled by Claude Sonnet 5.5 (claude-sonnet-5-5) for Andr\'as Kornai}
\date{}
\begin{document}
\maketitle

\section{Introduction}
Kolbila (ISO 639-3 klc) is a Central Adamawa (Samba-Duru) language of Cameroon, spoken along the Ngaound\'er\'e--Garoua road area. Speaker estimates vary (about 2,000 to 4,500). Very little material was found.

\section{Primary Educational Resources}
No primers, textbooks or school programmes were found.

\section{Linguistic and Reference Materials}
\begin{itemize}
\item \emph{Linguistische Beschreibung des Kolbila} (2016, German), the first scholarly monograph on the language, covering phonology and word classes: \url{https://www.koeppe.de/titel_print_linguistische-beschreibung-des-kolbila}
\item Wikipedia overview: \url{https://en.wikipedia.org/wiki/Kolbila_language}
\end{itemize}

\section{Religious and Community Resources}
None found.

\section{Digital and Audio Resources}
\begin{itemize}
\item Glosbe hosts crowd-sourced Kolbila word pairs (quality uncertain): \url{https://glosbe.com/klc/lib}
\end{itemize}

\section{Curriculum Development Context}
No published Kolbila dictionary was found. Curriculum work would need to start from the 2016 grammar and new community-based documentation and orthography work.

\end{document}